\section{Discussion}
\label{sec:discussion}

The results show a clear difference between the two commitment strategies. The na{\"\i}ve construction is useful as a simple baseline because it commits directly to the occupied ship coordinates, but it is limited by the arity of the $\poseidon$ hash. Since standard Battleship requires 17 occupied coordinates, and each coordinate consists of both an $x$ and a $y$ value, the direct $\poseidon$ commitment cannot represent a full standard board. The Merkle tree construction removes this restriction by committing to a root and proving membership only for the queried square. This makes the Merkle approach the more meaningful result for a realistic version of this game.

The timing results suggest that our prototype is dominated by fixed overheads rather than the cost of adding ship coordinates. The Merkle benchmark shows only small growth over the tested board sizes, which is consistent with the logarithmic Merkle path checked by the circuit. At the same time, however, the absolute running time remains large enough that the current implementation would feel slow in a real\-/time game. In practice, witness generation and snarkjs invocation likely contribute most to the user\-/visible latency. 

The most important cryptographic limitation of our game is the trusted setup required by $\groth$. Before either player can generate or verify proofs, the circuit must be processed into public proving and verification parameters. During this setup procedure, we create secret intermediate randomness. The security of the system depends on the assumption that we discard this randomness and it is not known to either player. If a malicious party retained it, they could produce proofs that verify even when the underlying Battleship statement is false. Thus, the protocol removes trust in a game server by introducing trust in the setup process. In our case, we assume that the game developer is trusted, but that does not preclude collusion between the developer and one of the players.

The implementation is also limited by timing. Although the Merkle construction scales better than the na{\"\i}ve construction, the current proving pipeline is still too slow to be invisible to the players. Since each move requires the defender to generate a proof before the attacker can accept, proof generation directly affects turn latency. If used on a game that was more real\-/time than Battleship, this could become a major barrier to deployment. Thus, although our results support Merkle tree design, they do not show this implementation could scale to a real\-/time game experience. A potential avenue for further research would be to prove a batch of previous game states at regular intervals, amortizing the proof generation cost but requiring that game states are retroactively checked.

The privacy guaranty is also narrower than what one might want from a fully private hidden\-/information game. Although our protocol hides the unrevealed board squares of the defender, the guesses of the attacker are public. This means that the defender may be able to infer the attacker's search strategy over time, so it does not protect the privacy of how a player chooses moves. A stronger formulation of the problem could hide both the defender's board and the attacker's targeting strategy.
